\begin{afterword}
\summaryhead

The post asks why people pursue their goals so badly, like an aspiring comedian who studies by watching old cartoons. The author's answer is that most possible actions are ineffective, just as most answers on a calculus test are wrong. Nothing in evolution or culture pushes us strongly toward the few that work.

People do have goals in a limited sense. They tell stories about them and feel good or bad about the results. But they do not automatically ask what they are trying to achieve, measure progress, gather information, try different methods, and put their effort where it works. The deeper reason, she suggests, is that the part of us that reasons in words is much better than the part that drives what we do. Knowing that exercise is good does not make us want to exercise. She ends by asking readers how they have trained themselves to be more strategic.

\responsehead

This is a strong post. The calculus-test analogy turns the question around: failure is the default and needs no special explanation, so success is what needs explaining. The list of habits is practical and has aged well. The glass floor captures the gap between knowing and wanting in one image.

Its main weakness is that it never counts the cost of being strategic. Searching, planning and testing take time and attention. Herbert Simon made this point in 1956 when he coined the word ``satisficing.'' Real people, with limited time and knowledge, stop at the first option that is good enough, and often they are right to. Some of the post's examples are real mistakes, such as choosing a career without comparing salaries. Others, such as two-finger typing or unplanned Saturdays, may be sensible trade-offs. The post needed a paragraph on how to tell the two apart. It also reaches for an unsupported number (5\% of people) and a grand claim (humans ``only just on the cusp of general intelligence'') that the argument does not need.

The post's picture of two systems, one that reasons in words and one that drives action, belongs to a family of ``dual-process'' theories. Keith Stanovich and Richard West named the two systems System 1 and System 2 in 2000. Daniel Kahneman made the terms famous in \textsc{Thinking, Fast and Slow} in 2011, a year after this post. The older name for knowing the better course and taking the worse is akrasia, which Aristotle discussed in the \textsc{Nicomachean Ethics}.

The post also had consequences. In 2012 its author co-founded the Center for Applied Rationality, which ran workshops teaching habits much like the ones on her list.

As an essay, it would earn high marks. The teacher would ask for the concrete examples to move from the footnote into the body, for the typos to be fixed, and for the invented statistic to go.
\end{afterword}
